\documentclass[11pt]{ctexart}
\usepackage[a4paper,margin=1in]{geometry}
\usepackage{graphicx}
\usepackage{xcolor}
\usepackage{hyperref}
\definecolor{refgray}{gray}{0.45}
\title{\textbf{Market Views｜全球投行叙事汇编}\\[0.4em]{\large 增长韧性叠加通胀退潮，AI基础设施扩张从硬件向软件价值重估，美国电力需求结构性拐点确认}}
\author{KC桌面 自动生成}
\date{260714}
\begin{document}
\maketitle
\tableofcontents
\newpage
\section{一页摘要}
\begin{itemize}
  \item 全球宏观呈现增长韧性（企业盈利上修中位数约8\%）与通胀退潮（中东冲击二季度见顶、中国PPI或6月见顶）并存格局，但政策路径分化：韩国被迫加息，中国央行维持观察，日本受益于利率正常化。
  \item AI资本开支持续扩张，北美五大超大规模云厂商2028年合计资本开支预计达1.4万亿美元，算力四年四倍；但市场对WFE规模预期（2000亿美元）已高于GS基准（1870亿美元），存在修正空间。
  \item AI从训练转向推理驱动服务器CPU需求结构变化，智能体AI服务器CPU以155\%年复合增速成为最大增量，2028年全球服务器CPU出货量预计达6800万颗。
  \item 中国AI硬件（科创板50）相对软科技（恒生科技）超额收益达68个百分点创历史纪录，但GS认为轮动需等待互联网板块盈利拐点（预计2Q-3Q26），而非估值修复。
  \item 美国电力需求增速将从历史低位（0.35\% CAGR）跃升至约3\% CAGR（2025-2030年），数据中心是最大增量来源，其电力需求占比将从6\%升至10\%；Bernstein预测在主流机构中偏保守，分歧在于建设进度和电气化渗透率。
  \item 天然气需求增长由LNG出口、电力用气和供给侧损耗三支柱驱动，到2030年总需求预计增长约24\%至154 bcfd，亨利港价格中周期预期为5美元/百万英热单位。
  \item 亚洲奢侈品市场（日本、韩国、香港）在2026年二季度表现强劲，硬奢持续优于软奢，AI相关财富效应和游客消费是主要驱动力；汽车行业在中国市场深度调整，但欧美支撑和成本优化使利润底线暂稳。
  \item 代币化RWA市场正从发行实验转向效用落地，Robinhood Chain等实际部署验证了用户对代币化资产流动性的需求；中国NPBFI工具是当前双重约束下的理性准财政替代方案，GS估算拉动GDP约0.5个百分点，但效果集中在2026年底至2027年初。
\end{itemize}
\section{全球投行叙事汇编}
今日纳入 30 篇投行/券商研究，按 5 个市场、资产与产业主题系统整理。
\newpage
\subsection{宏观与利率：增长韧性、通胀退潮与政策分化}
\textbf{全球宏观呈现增长韧性（企业盈利上修）与通胀退潮（中东冲击减弱、中国PPI见顶）并存的格局，但政策路径分化：韩国被迫加息，中国央行维持观察，日本受益于利率正常化，人民币中间价模型暗示升值空间。}
\subsubsection*{主流共识}
\begin{itemize}
  \item 全球企业盈利预期在二季度持续上修，中位数增长约8\%，反映基本面真实改善而非分析师宽松调整。
  \item 中东冲突对通胀的增量影响已在二季度见顶，三季度将明显下降，除非油价再度飙升。
  \item 中国PPI通胀可能在6月见顶，上游价格压力缓解，企业成本端最紧张阶段已过。
\end{itemize}
\subsubsection*{分歧与条件差异}
\begin{itemize}
  \item 韩国AI出口红利能否转化为国内增长：GS认为关键在汇率和财政政策，韩元逆势调整正在抵消贸易条件改善的溢出效应。
  \item 中国央行政策方向：市场部分预期宽松，但GS判断2026年不会再有政策利率或准备金率下调，央行更倾向结构性工具。
  \item 日本银行体系盈利前景：GS上调三大银行目标价，认为市场定价仍停留在旧框架，但三菱日联的海外利率敏感度构成风险。
\end{itemize}
\subsubsection*{机构观点}
\paragraph{GS} 韩国AI红利传导受阻于韩元逆势调整。贸易条件改善约20\%本应推升汇率，但韩元因海外投资和再平衡操作走弱，削弱了出口收入向国内购买力的转化。韩元每贬值5\%推高通胀约30bp，可能迫使韩国央行提前加息。
{\small\color{refgray}数据：韩国Q2贸易条件同比改善近20\%\par}
{\small\color{refgray}数据：韩元对美元每贬值5\%推高通胀约30bp\par}
{\small\color{refgray}数据：GS预计韩国央行7月16日启动加息\par}
{\small\color{refgray}边际变化：此前市场认为AI出口红利将自然传导至国内，但GS指出汇率渠道失效，政策干预成为关键变量。\par}
\paragraph{GS} 中国6月通胀数据确认PPI见顶，央行维持观察窗口。PPI同比4.1\%但环比明显走弱，GS下调下半年PPI预测至2026年2\%、2027年0\%。央行二季度例会承认“供给强、需求弱”，但未释放宽松信号，预计2026年无降息降准。
{\small\color{refgray}数据：中国6月CPI同比1.0\%，PPI同比4.1\%\par}
{\small\color{refgray}数据：GS预测2026年全年PPI为2\%，2027年为0\%\par}
{\small\color{refgray}数据：央行二季度例会未暗示全面宽松\par}
{\small\color{refgray}边际变化：市场此前担忧PPI持续上行，但6月环比转弱证伪该预期，政策焦点从宽松转向结构性工具。\par}
\paragraph{GS} 日本银行业受益于三重推力：贷款增长（大型银行同比8.4\%）、日央行两次加息扩大净息差、资本市场活跃。GS上调三菱日联目标价15\%至3900日元，三井住友上调14\%至7300日元，预计ROE从10-12\%升至12-13\%。
{\small\color{refgray}数据：日本4-6月大型银行贷款同比+8.4\%\par}
{\small\color{refgray}数据：GS假设日央行政策利率2027财年达0.75\%\par}
{\small\color{refgray}数据：三菱日联CET1比率9.2\%，低于目标区间9.5-10.5\%\par}
{\small\color{refgray}边际变化：市场将银行股视为利率滞后指标，但GS认为利率调整已在资金合同和债券组合中兑现，盈利上修有实物证据支撑。\par}
\paragraph{GS} 中东冲突对通胀的增量影响已过峰值。油价较4月底峰值低30\%，汽油价格跌15\%，航空燃油跌35\%。模型显示大宗商品向消费价格的传导在Q2见顶，Q3-Q4递减。核心PCE 6月环比约24bp，之后维持在20-23bp。
{\small\color{refgray}数据：油价较4月底峰值低约30\%\par}
{\small\color{refgray}数据：汽油价格下跌15\%，航空燃油下跌35\%\par}
{\small\color{refgray}数据：核心PCE 6月环比预计24bp，之后20-23bp\par}
{\small\color{refgray}边际变化：此前市场担忧冲突持续推高通胀，但GS指出增量影响已在Q2见顶，除非油价重回100美元/桶，否则通胀风险下降。\par}
\paragraph{JPM} 二季度盈利预期走高并非陷阱，而是基本面改善的真实映射。美国22\%、欧元区12\%的市值加权增速受IT和能源扭曲，中位数约8\%更可达成。欧元区EPS修正已转正，与美国差距接近弥合，信贷增长和OECD领先指标提供支撑。
{\small\color{refgray}数据：美国Q2 EPS中位数增长预期约8\%\par}
{\small\color{refgray}数据：欧元区Q2 EPS中位数增长预期约8\%\par}
{\small\color{refgray}数据：MSCI AC World多数行业录得净上调\par}
{\small\color{refgray}边际变化：通常财报季前分析师会下调预期，但本次2026年EPS预测持续上修，且范围广泛，反映真实宏观改善。\par}
\paragraph{NOM} 人民币中间价模型投影6.7841，较前次下移148bp，计入逆周期因子后为6.7902。模型误差收窄暗示逆周期因子使用减少，定价机制回归常态。下半年APEC、中美会晤等事件构成升值窗口。
{\small\color{refgray}数据：无逆周期因子模型投影6.7841\par}
{\small\color{refgray}数据：较前次预测6.7989下移148bp\par}
{\small\color{refgray}数据：计入逆周期因子后投影6.7902\par}
{\small\color{refgray}边际变化：模型误差持续收窄，说明官方减少对中间价的人为干预，市场供求在定价中作用上升。\par}
\subsubsection*{关键数据}
\begin{itemize}
  \item 韩国Q2贸易条件同比改善近20\%，但韩元逆势调整
  \item 中国6月PPI同比4.1\%但环比走弱，GS预测2026年PPI为2\%
  \item 日本4-6月大型银行贷款同比+8.4\%
  \item 油价较4月底峰值低30\%，汽油价格跌15\%
  \item 美国Q2 EPS中位数增长预期约8\%
  \item 人民币中间价模型投影6.7841，较前次下移148bp
  \item GS预计韩国央行7月16日启动加息
  \item 中东冲突对核心PCE的增量影响在Q2见顶
\end{itemize}
\begin{figure}[htbp]
\centering
\includegraphics[width=0.88\linewidth]{figures/fig_034.jpg}
\caption{Exhibit 1 - Exhibit 2: Korea's overall terms of trade have improved sharply, pushing up its GDP deflator}
\end{figure}
\begin{figure}[htbp]
\centering
\includegraphics[width=0.88\linewidth]{figures/fig_072.jpg}
\caption{Exhibit 1 - Exhibit 1: Oil Prices Have Risen Only Moderately During the Latest Attacks and Remain Roughly 30\% Below Their Wartime Peak; Gasoline Prices Are Down 15\%, and Jet Fuel Is Down 35\% The prices of other Persian Gulf exports such as}
\end{figure}
\begin{figure}[htbp]
\centering
\includegraphics[width=0.88\linewidth]{figures/fig_077.jpg}
\caption{Figure 1 - Figure 1: S\&P500 and MSCI Eurozone 2026e EPS levels In our bullish 2026 Year Ahead view, a strong earnings outlook was one of the main supports. This went against the view of many investors who thought consensus earnings expecta}
\end{figure}
\begin{figure}[htbp]
\centering
\includegraphics[width=0.88\linewidth]{figures/fig_078.jpg}
\caption{Figure 2 - Figure 3: MSCI AC World ytd}
\end{figure}
{\small\color{refgray}本节综合 6 篇研究；涉及 GS、JPM、NOM。}
\newpage
\subsection{AI基础设施扩张：从芯片到系统，从硬件到软件的结构性重塑}
\textbf{AI基础设施正经历从训练到推理、从单一芯片到系统级架构、从硬件主导到软硬协同的结构性转变。半导体设备、EDA工具、服务器CPU和AI平台公司均受益于这一趋势，但各环节的定价和预期分化显著，市场对硬件端的乐观已部分透支，而软件和AI平台的价值重估刚刚开始。}
\subsubsection*{主流共识}
\begin{itemize}
  \item AI资本开支持续扩张，北美五大超大规模云厂商2028年合计资本开支预计达1.4万亿美元，算力从2025年30GW增至120GW，四年四倍。
  \item 半导体设备需求与AI基础设施扩张高度相关，前端设备（东京电子、SCREEN等）获得市场最多关注，WFE市场2027年预期在1870-2000亿美元区间。
  \item AI从训练转向推理驱动服务器CPU需求结构变化，智能体AI服务器CPU以155\%年复合增速成为最大增量，2028年全球服务器CPU出货量预计达6800万颗。
\end{itemize}
\subsubsection*{分歧与条件差异}
\begin{itemize}
  \item 市场对WFE规模的预期（2000亿美元）已高于GS基准预测（1870亿美元），若存储厂资本开支节奏或设备涨价落地不及预期，当前定价存在修正空间。
  \item 中国AI硬件（科创板50）相对软科技（恒生科技）超额收益达68个百分点，创历史纪录，但GS认为轮动需等待互联网板块盈利拐点（预计2Q-3Q26），而非估值修复。
  \item Meta的AI能力对外输出（Muse Spark 1.1、API定价为竞品1/4）被JPM和MS视为未被定价的期权，但市场对其资本开支扩张的投入产出比仍存疑。
  \item 英伟达的CPU业务（Vera独立机架）和ASIC渗透（某前沿模型份额从0升至50\%）显示其增长边界超出市场想象，但MS强调内存短缺将持续数年，需通过计算-网络-内存协同优化。
\end{itemize}
\subsubsection*{机构观点}
\paragraph{Bernstein} 存储与半导体设备的相关性被市场高估，2011-2026年两者12个月滚动相关性均值仅0.54，远低于设备与整体半导体板块的0.82。在2015-2016年存储下跌16.2\%时设备上涨21.9\%，2021-2022年存储下跌34\%时设备上涨15.3\%，设备板块可走出独立行情。长期看，存储2019-2025年累计涨幅661\%领先设备，但截至2026年2月两者15年涨幅均达约36倍，均值回归下设备存在空间。
{\small\color{refgray}数据：存储与设备12个月滚动相关性均值0.54（2011-2026）\par}
{\small\color{refgray}数据：2015-2016年设备上涨21.9\% vs 存储下跌16.2\%\par}
{\small\color{refgray}数据：2021-2022年设备上涨15.3\% vs 存储下跌34\%\par}
{\small\color{refgray}数据：存储2019-2025年累计涨幅661\%\par}
{\small\color{refgray}边际变化：市场此前普遍假设存储与设备高度绑定，Bernstein通过七个周期阶段数据证明设备可独立于存储表现，且当前存储相对设备的溢价已高于历史均值。\par}
\paragraph{GS} EDA行业正从AI浪潮旁观者变为核心受益者。全球芯片设计工程师到2030年短缺约7.2万名，而AI设计工具可等效替代约4.8万名工程师。GS估算这为EDA行业带来每年约37亿美元增量收入，当前市场共识几乎为零。Cadence营收增速可能从13\%提升至25\%以上，Synopsys可达20\%左右，增量贡献最早2026年下半年显现。
{\small\color{refgray}数据：2030年全球芯片设计师短缺约7.2万名\par}
{\small\color{refgray}数据：AI工具可等效替代约4.8万名工程师\par}
{\small\color{refgray}数据：EDA行业每年增量收入约37亿美元\par}
{\small\color{refgray}数据：Cadence营收增速可能从13\%提升至25\%以上\par}
{\small\color{refgray}边际变化：市场此前对EDA公司的预期停留在12\%-15\%稳定增长，GS认为AI定制芯片需求加剧工程师短缺，AI工具将成为填补缺口的关键手段，增量收入未被定价。\par}
\paragraph{GS} 市场对日本半导体前端设备的预期已超前。7月初香港路演中，投资者对2027年WFE市场预期集中在2000亿美元，高于GS基准1870亿美元。前端设备（东京电子、SCREEN、Kokusai Electric）获最多乐观情绪，设备涨价讨论升温——东京电子已推动涨价，市场预期扩散。后端设备（Disco、Advantest）关注度相对冷清。
{\small\color{refgray}数据：投资者对2027年WFE预期约2000亿美元\par}
{\small\color{refgray}数据：GS基准预测1870亿美元\par}
{\small\color{refgray}数据：前端设备获最多乐观情绪，后端相对冷清\par}
{\small\color{refgray}数据：东京电子已推动涨价，市场预期扩散\par}
{\small\color{refgray}边际变化：市场不仅接受更高WFE总量假设，还在为设备商利润率改善提前定价，但GS认为若存储厂资本开支节奏或涨价落地不及预期，当前定价存在修正空间。\par}
\paragraph{GS} 工业富联的核心竞争力不在组装，而在为全球高质量客户完成最新技术产品从0到1的量产导入。越南工厂超15年经验、自研自动化产线快速复制能力，以及客户结构从大型云厂商向NeoCloud等新兴算力需求方扩展，使其在AI基础设施扩张周期中获得议价权。GS上调2026-2028年全球AI服务器隐含芯片出货量预期14\%-22\%，中国云服务商资本开支上调37\%-55\%。
{\small\color{refgray}数据：越南工厂超15年生产经验\par}
{\small\color{refgray}数据：2026-2028年全球AI芯片出货量预期上调14\%-22\%\par}
{\small\color{refgray}数据：中国云服务商资本开支上调37\%-55\%\par}
{\small\color{refgray}数据：越南工厂是海外基地中利润率最高的\par}
{\small\color{refgray}边际变化：市场此前将工业富联视为普通组装厂，GS通过越南工厂实地调研揭示其自动化方案商定位，以及客户愿意为产能多元化支付溢价的结构性窗口。\par}
\paragraph{GS} 中国AI硬件（科创板50）年内跑赢恒生科技68个百分点，创历史纪录，但AI整体并非泡沫——AI通过效率提升和新利润创造带来的潜在宏观收益比当前定价高出50\%-100\%。轮动至H股软科技的必要条件是盈利可见修复，而非估值便宜。MSCI中国一季度利润同比下降8\%，互联网板块累计补贴亏损超1800亿元，预计2Q-3Q26出现经营利润拐点。
{\small\color{refgray}数据：科创板50相对恒生科技超额收益68个百分点\par}
{\small\color{refgray}数据：AI潜在宏观收益比当前定价高50\%-100\%\par}
{\small\color{refgray}数据：MSCI中国一季度利润同比下降8\%\par}
{\small\color{refgray}数据：互联网板块累计补贴亏损超1800亿元\par}
{\small\color{refgray}边际变化：市场关注A股硬件与港股软件的分化，GS明确指出轮动条件不是估值便宜，而是盈利可见修复，且当前估值已反映较悲观情景。\par}
\paragraph{GS} 东南亚AI供应链迁移不是成本套利，而是客户意愿支撑的结构性产能再布局。泰国/越南工厂综合生产成本比中国大陆高15\%-20\%，但客户愿意通过预付或补贴设备锁定产能。越南工厂原材料仍需从中国进口，但客户对供应安全的重视超过成本优化。中国大陆工人产出效率是泰国的3倍，但自动化正在缩小差距。
{\small\color{refgray}数据：东南亚工厂综合成本比中国大陆高15\%-20\%\par}
{\small\color{refgray}数据：中国大陆工人产出是泰国工人的3倍\par}
{\small\color{refgray}数据：越南工厂原材料（如覆铜板）仍需从中国进口\par}
{\small\color{refgray}数据：所有被访企业扩产方向全部指向AI基础设施\par}
{\small\color{refgray}边际变化：市场此前认为东南亚迁移是成本驱动，GS调研发现客户主动要求并愿意付费，且东南亚正形成围绕AI需求的独立生态系统。\par}
\paragraph{JPM} Meta的AI策略从自用转向对外变现，Muse Spark 1.1模型在智能体和编码能力上显著提升，API定价为头部模型的25\%，瞄准1500万广告主。下一代模型Watermelon性能可能匹配GPT 5.5。JPM认为这些变化为盈利提供上行空间，但需关注模型能力追赶速度能否支撑从广告公司升级为AI平台公司。
{\small\color{refgray}数据：Meta股价从6月底低点反弹超23\%\par}
{\small\color{refgray}数据：Muse Spark 1.1 API定价为头部模型的25\%\par}
{\small\color{refgray}数据：下一代模型Watermelon可能匹配GPT 5.5\par}
{\small\color{refgray}数据：Meta 2026年资本支出预计1420亿美元（+104\%）\par}
{\small\color{refgray}边际变化：市场此前认为Meta在AI模型竞赛中落后，JPM指出Muse Spark 1.1缩小了与Anthropic、OpenAI、Google的差距，且Meta首次将AI能力对外输出。\par}
\paragraph{JPM} 服务器CPU市场正经历结构性变化，AI推理与智能体工作负载成为主要驱动力。2028年全球服务器CPU出货量将从2025年2600万颗增至6800万颗（38\%年复合增长），其中智能体AI服务器CPU以155\%年复合增速成为最大增量。加速器与CPU配比从4:1压缩至2:1甚至1:1，ARM架构份额从2025年22\%升至2028年43\%。
{\small\color{refgray}数据：2028年服务器CPU出货量6800万颗（2025年2600万颗）\par}
{\small\color{refgray}数据：智能体AI服务器CPU年复合增长155\%\par}
{\small\color{refgray}数据：传统通用服务器CPU年复合增长约5\%\par}
{\small\color{refgray}数据：加速器与CPU配比从4:1降至2:1甚至1:1\par}
{\small\color{refgray}边际变化：市场此前关注GPU，JPM指出AI从训练转向推理正在重新定义CPU需求结构，智能体AI CPU是突破性增长向量，且ARM架构加速渗透。\par}
\paragraph{MS} 英伟达增长正在加速，三个引擎同时展开：AI实验室（某前沿模型从ASIC切回英伟达，份额从0升至50\%）、超大规模云（网络+CPU扩大SAM）、主权/企业（竞争最弱）。CPU业务200亿美元目标中近一半来自独立机架，Vera芯片下半年专为单线程AI任务设计。管理层主动向价值型投资者敞开大门，现金流50\%以上回馈股东。
{\small\color{refgray}数据：某前沿模型从ASIC切回英伟达，份额从0升至50\%\par}
{\small\color{refgray}数据：CPU业务目标200亿美元，近一半来自独立机架\par}
{\small\color{refgray}数据：Vera芯片下半年专为单线程AI任务设计\par}
{\small\color{refgray}数据：现金流50\%以上回馈股东\par}
{\small\color{refgray}边际变化：市场认为英伟达增长接近天花板，MS通过NDR揭示三个增长引擎边界更宽，且公司正从成长股向价值股定位转变。\par}
\paragraph{MS} 北美五大超大规模云厂商2027/2028年资本开支上调至1.2万亿和1.4万亿美元，算力从2025年30GW增至2028年120GW（四年四倍）。从破土动工到数据中心投产周期拉长至三年，迫使厂商提前建设。每GW建造成本上调约20\%，内存成本从低个位数占比升至高双位数。META被列为首选，五个未被定价的盈利期权合计可能带来每股约10美元额外收益。
{\small\color{refgray}数据：2028年五大云厂商资本开支1.4万亿美元\par}
{\small\color{refgray}数据：算力从30GW增至120GW（四年四倍）\par}
{\small\color{refgray}数据：每GW建造成本上调约20\%\par}
{\small\color{refgray}数据：META五个盈利期权合计每股约10美元额外收益\par}
{\small\color{refgray}边际变化：市场对AI资本开支持续性存疑，MS给出更激进预测，并指出META的API业务和AI广告工具等期权未被市场定价。\par}
\subsubsection*{关键数据}
\begin{itemize}
  \item 存储与设备12个月滚动相关性均值0.54（2011-2026）
  \item EDA行业每年增量收入约37亿美元（GS估算）
  \item 投资者对2027年WFE预期约2000亿美元 vs GS基准1870亿美元
  \item 科创板50相对恒生科技超额收益68个百分点
  \item 2028年全球服务器CPU出货量6800万颗（2025年2600万颗）
  \item 智能体AI服务器CPU年复合增长155\%
  \item 北美五大云厂商2028年资本开支1.4万亿美元
  \item 算力从30GW增至120GW（四年四倍）
  \item META五个盈利期权合计每股约10美元额外收益
\end{itemize}
\begin{figure}[htbp]
\centering
\includegraphics[width=0.88\linewidth]{figures/fig_019.jpg}
\caption{Exhibit 1 - EXHIBIT 1: Correlation between memory and other parts of semis has historically been lower and more volatile. As we established above that semicap and the broader SOX is tightly correlated, we shift our attention to how memory}
\end{figure}
\begin{figure}[htbp]
\centering
\includegraphics[width=0.88\linewidth]{figures/fig_044.jpg}
\caption{Exhibit 1 - \#\# Chip design labor shortage could drive a material EDA revenue opportunity through 2030 Bottom Line: The semiconductor industry has a long-standing labor problem, with the world producing new chip designers at just 2\%-3\% a year—yielding \textbackslash{}\textasciitilde{}255k engineers by 2}
\end{figure}
\begin{figure}[htbp]
\centering
\includegraphics[width=0.88\linewidth]{figures/fig_057.jpg}
\caption{Exhibit 1 - Exhibit 1: The AI bottleneck, Hard vs. Software, or AI supplier vs. AI capex spender trade has been a dominant equity theme globally Exhibit 2: Indexes with higher AI Hardware exposures have outperformed ytd in the Chinese equity}
\end{figure}
\begin{figure}[htbp]
\centering
\includegraphics[width=0.88\linewidth]{figures/fig_082.jpg}
\caption{Figure 1 - Figure 1: Server CPU shipment trend and CAGR Figure 2: Server CPU revenue TAM forecast and CAGR \#\# Top-down methodology}
\end{figure}
{\small\color{refgray}本节综合 10 篇研究；涉及 Bernstein、GS、JPM、MS。}
\newpage
\subsection{美国电力与天然气：结构性需求拐点与供给竞争}
\textbf{美国电力需求正经历从长期停滞（0.35\% CAGR）到结构性加速（约3\% CAGR）的根本性转变，数据中心、制造业回流和电气化三重力量驱动；天然气作为发电主力燃料，受益于电力需求增长和LNG出口扩张，到2030年总需求有望增长约24\%。储能电池成为全球能源转型的弹性变量，竞争格局重塑中。}
\subsubsection*{主流共识}
\begin{itemize}
  \item 美国电力需求增速将从历史低位（0.35\% CAGR）跃升至约3\% CAGR（2025-2030年），数据中心是最大增量来源，其电力需求占比将从6\%升至10\%。
  \item 天然气需求增长由LNG出口、电力用气和供给侧损耗三支柱驱动，到2030年总需求预计增长约24\%至154 bcfd，亨利港价格中周期预期为5美元/百万英热单位。
  \item 储能电池（ESS）是未来电池行业增长的最大弹性变量，竞争格局尚未固化，产能利用率将加速分化。
\end{itemize}
\subsubsection*{分歧与条件差异}
\begin{itemize}
  \item Bernstein对电力需求增速的预测（约3\% CAGR）在主流机构中偏保守，IEA、EIA、Rystad预测区间为2\%-5\%，分歧在于数据中心建设进度和电气化渗透率。
  \item 天然气发电用气2025年尚未明显增长（因煤电替代效应），但Bernstein预计2026年起每年新增约2 bcfd；市场对Haynesville钻机回升的担忧可能高估，因回升主要来自低产井运营商。
  \item 宁德时代向集成储能解决方案提供商的转型是否被充分定价：GS认为市场低估，储能业务估值倍数（25x）高于动力电池（20x），但执行和贸易政策风险仍存。
\end{itemize}
\subsubsection*{机构观点}
\paragraph{Bernstein} 美国电力需求正经历根本性变化，预计2025-2030年CAGR约3\%，远高于2000-2024年的0.35\%。数据中心是最大增量，当前占电力需求6\%，2030年升至10\%；制造业回流、电动车和热泵普及贡献额外约170 TWh。天然气需求同步扩张，到2030年总需求增长约24\%至154 bcfd，其中LNG出口从18 bcfd增至33 bcfd，电力用气累计增加约6 bcfd（2026年起每年约2 bcfd）。供给端Permian和Appalachia钻机数自2024年初分别下降25\%和21\%，Haynesville回升有限，供给纪律明显。
{\small\color{refgray}数据：美国电力需求CAGR 3\%（2025-2030年），对比2000-2024年0.35\%\par}
{\small\color{refgray}数据：数据中心装机从45 GW增至86 GW，新增电力需求约230 TWh\par}
{\small\color{refgray}数据：PUE假设1.1x-1.5x导致数据中心需求区间200-270 TWh\par}
{\small\color{refgray}数据：天然气总需求增长24\%至154 bcfd（2030年）\par}
{\small\color{refgray}边际变化：从过去二十年需求停滞转向结构性加速，数据中心和电气化成为新驱动力；天然气需求增长由LNG和电力双引擎驱动，供给纪律强化。\par}
\paragraph{GS} 宁德时代增长驱动力正从动力电池切换至储能电池（ESS），市场对其向集成储能解决方案提供商的转型定价不足。中国电池总出货量到2030年弹性较高，储能是最大变量。宁德时代在储能领域的份额有望回升至与动力电池相当的水平，其储能电池出货结构将向集成解决方案倾斜，带动单位毛利提升。GS对动力电池和储能业务分别估值，储能估值倍数更高（25x vs 20x）。
{\small\color{refgray}数据：中国电池总出货量到2030年弹性较高，储能是最大变量\par}
{\small\color{refgray}数据：储能电池竞争格局重塑，产能利用率分化加速\par}
{\small\color{refgray}数据：宁德时代储能电池份额有望回升至与动力电池相当\par}
{\small\color{refgray}数据：储能业务估值倍数25x，高于动力电池的20x\par}
{\small\color{refgray}边际变化：从电池电芯供应商转向集成储能解决方案提供商，储能业务估值逻辑独立化且倍数更高。\par}
\subsubsection*{关键数据}
\begin{itemize}
  \item 美国电力需求CAGR 3\%（2025-2030年），对比2000-2024年0.35\%
  \item 数据中心装机从45 GW增至86 GW，新增电力需求约230 TWh
  \item 天然气总需求增长24\%至154 bcfd（2030年）
  \item LNG液化能力从18 bcfd（2026年初）扩至33 bcfd（2030年底）
  \item 电力用气累计增加约6 bcfd（2025-2030年）
  \item Permian和Appalachia钻机数分别下降25\%和21\%
  \item 亨利港中周期价格预期5美元/百万英热单位
  \item 储能电池竞争格局重塑，产能利用率分化加速
  \item 宁德时代储能业务估值倍数25x，高于动力电池的20x
\end{itemize}
\begin{figure}[htbp]
\centering
\includegraphics[width=0.88\linewidth]{figures/fig_001.jpg}
\caption{EXHIBIT 1 - EXHIBIT 2: Our in-house model allows us to sensitize demand and supply on various fronts}
\end{figure}
\begin{figure}[htbp]
\centering
\includegraphics[width=0.88\linewidth]{figures/fig_011.jpg}
\caption{EXHIBIT 1 - EXHIBIT 1: Three pillars grow gas demand by 24\% by 2030 Growing to 2030 gas demand (bcfd) As we update our US power model, power demand growth provides significant growth in gas demand used in generation. Total power demand (TWh)}
\end{figure}
\begin{figure}[htbp]
\centering
\includegraphics[width=0.88\linewidth]{figures/fig_002.jpg}
\caption{EXHIBIT 4 - EXHIBIT 4: Our scenarios for potential data center power demand range from +200TWh to +270TWh reflecting a PUE range of 1.1x to 1.5x U.S. data center electricity demand (TWh) We see commercial demand overall growing at a 3.6\% CAG}
\end{figure}
\begin{figure}[htbp]
\centering
\includegraphics[width=0.88\linewidth]{figures/fig_012.jpg}
\caption{EXHIBIT 2 - EXHIBIT 3: ...with gas assuming a large piece of the pie EXHIBIT 4: We expect gas demand for power to grow at a 2.5\% CAGR through 2030 Implied Gas Demand for Power (bcfd)}
\end{figure}
{\small\color{refgray}本节综合 3 篇研究；涉及 Bernstein、GS。}
\newpage
\subsection{消费与汽车：分化加剧，利润修复先于销量复苏}
\textbf{全球消费与汽车板块呈现显著分化：高端消费受AI财富效应和区域旅游驱动，亚洲市场韧性超预期；汽车行业则面临中国市场深度调整，但欧美支撑和成本优化使利润底线暂稳。现制饮品竞争格局趋于有序，耐克等品牌正经历主动修复，利润拐点可能先于收入到来。}
\subsubsection*{主流共识}
\begin{itemize}
  \item 亚洲奢侈品市场（日本、韩国、香港）在2026年二季度表现强劲，硬奢（珠宝腕表）持续优于软奢（皮具成衣），AI相关财富效应和游客消费是主要驱动力。
  \item 汽车行业在中国市场面临深度调整，宝马中国区销量同比下滑约30\%，保时捷FY27全球交付量预期下调，但欧美市场增长和成本削减措施支撑利润指引。
  \item 现制饮品赛道竞争格局趋于有序，茶饮品牌进军咖啡规模有限，定价促销保持理性，头部品牌同店销售表现优于市场悲观预期。
\end{itemize}
\subsubsection*{分歧与条件差异}
\begin{itemize}
  \item 对耐克修复节奏的分歧：Bernstein认为FY27仍是转型年，Lifestyle收入中个位数下滑，利润拐点先于收入；市场部分观点期待更快复苏。
  \item 对宝马中国前景的分歧：GS下调中国合资企业利润率预测至1.3\%，但管理层维持全年汽车业务利润率1\%-3\%指引，欧美韧性能否持续存疑。
  \item 对保时捷利润驱动力的分歧：GS强调911组合定价和SG\&A削减足以抵消销量下滑，市场可能低估成本优化力度。
  \item 对B站货币化潜力的分歧：GS认为社区生态已进化为货币化飞轮，AI加速广告增长；部分投资者仍聚焦游戏管线不确定性。
\end{itemize}
\subsubsection*{机构观点}
\paragraph{Bernstein} 耐克正经历主动品牌修复，FY27仍是转型年。Lifestyle（含Jordan）占品牌收入约60\%，全年仍将录得中个位数下滑，经典款（AF1、AJ1、Dunk）对总销售影响从FY26的4个百分点收窄至约2个百分点。但利润率可能先于收入见底：Q1起毛利率逐季改善，主因减少折扣和清理渠道库存。EPS在FY27可基本持平，支撑来自SG\&A严格管控而非收入增长。
{\small\color{refgray}数据：Lifestyle占Nike品牌收入约60\%\par}
{\small\color{refgray}数据：经典款对FY26总销售影响约4个百分点，FY27收窄至约2个百分点\par}
{\small\color{refgray}数据：FY27 EPS可基本持平，支撑来自SG\&A管控\par}
{\small\color{refgray}边际变化：此前市场关注收入拐点，Bernstein强调利润拐点先于收入，成本纪律而非经营杠杆是短期EPS支撑。\par}
\paragraph{Bernstein} 全球奢侈品需求受AI相关财富效应带动，亚洲市场表现超预期。日本东京百货珠宝销售在2Q26仍保持20\%以上同比增长，韩国百货店4-5月增幅超30\%，香港珠宝销售同比增超20\%。硬奢（珠宝腕表）持续优于软奢，但差距收窄。若历史相关性有效，LVMH时装皮具2Q26有机销售增速可达+3\%，高于市场预期的+1.7\%。
{\small\color{refgray}数据：东京百货珠宝销售2Q26同比增20\%以上\par}
{\small\color{refgray}数据：韩国百货店4-5月奢侈品销售增幅超30\%\par}
{\small\color{refgray}数据：LVMH时装皮具2Q26有机销售增速预期+3\%，市场共识+1.7\%\par}
{\small\color{refgray}边际变化：市场此前关注中国消费疲软，Bernstein指出日本、韩国、香港的AI财富效应和游客红利是超预期来源。\par}
\paragraph{GS} B站社区生态已从护城河进化为货币化飞轮，AI成为加速器。广告增长来自三重引擎：流量增长（DAU同比+8\%，总使用时长+19\%）、用户消费力提升（平均年龄27岁，55\%在一二线城市）、深度场景货币化（搜索广告收入同比+100\%，播放页广告+80\%）。游戏管线目标未来2-3年新增至少2款长生命周期品类头部游戏。
{\small\color{refgray}数据：DAU同比+8\%，人均使用时长增加11分钟\par}
{\small\color{refgray}数据：搜索广告收入同比+100\%，播放页广告+80\%\par}
{\small\color{refgray}数据：2009年上传视频过去一年仍获203万次新增播放\par}
{\small\color{refgray}边际变化：市场此前聚焦游戏管线，GS强调AI定位为重大投入方向，且广告增长从流量、用户消费力、深度场景三个维度同步加速。\par}
\paragraph{GS} 宝马2Q26零售销量59.1万辆，同比-4.9\%，降幅几乎由中国市场驱动（GS估算中国区销量同比-30\%）。但欧洲和美国分别录得8\%和10\%增长，部分抵消中国调整。汽车业务EBIT利润率预估2.9\%，处于全年1\%-3\%指引区间内。中国合资企业全年利润率预测从2.3\%下调至1.3\%，但管理层维持全年指引不变。
{\small\color{refgray}数据：2Q26零售销量59.1万辆，同比-4.9\%\par}
{\small\color{refgray}数据：中国区销量同比-30\%，欧洲+8\%，美国+10\%\par}
{\small\color{refgray}数据：汽车业务EBIT利润率预估2.9\%\par}
{\small\color{refgray}数据：中国合资企业利润率预测从2.3\%下调至1.3\%\par}
{\small\color{refgray}边际变化：此前市场预期中国区销量下滑-24\%，实际-30\%超预期，但欧美增长和成本削减使利润指引维持不变。\par}
\paragraph{GS} 现制饮品竞争格局走向有序，茶饮品牌进军咖啡规模有限。古茗咖啡GMV预计2026年约80亿元，仅为瑞幸的低双位数百分比。定价促销保持理性：古茗每周仅一张9.9元咖啡券，瑞幸9.9活动仅限部分SKU。茶百道6月同店销售仅下滑低个位数，古茗同样维持低个位数下滑，市场悲观预期已有所调整。
{\small\color{refgray}数据：古茗2026年咖啡GMV预计约80亿元，为瑞幸低双位数百分比\par}
{\small\color{refgray}数据：茶百道6月同店销售下滑低个位数\par}
{\small\color{refgray}数据：古茗6月同店销售维持低个位数下滑\par}
{\small\color{refgray}边际变化：市场担忧茶饮品牌进军咖啡加剧竞争，GS指出规模差距和门店区位错位使竞争影响被高估，定价促销保持理性。\par}
\paragraph{GS} 保时捷2Q26利润修复超预期，核心驱动力并非销量。预计当季EBIT 7.85亿欧元，利润率8.5\%，显著高于市场共识6.6\%。911车型组合持续上移（ASP预计FY27同比+8\%），叠加SG\&A加速削减，足以抵消FY27全球交付量预期下调（GS预测同比-11\%）。第二轮人员优化计划接近敲定，涉及约4000名员工，规模可能超过第一轮。
{\small\color{refgray}数据：2Q26 EBIT预估7.85亿欧元，利润率8.5\%，市场共识6.6\%\par}
{\small\color{refgray}数据：911车型ASP FY27预计同比+8\%\par}
{\small\color{refgray}数据：FY27全球交付量预期同比-11\%\par}
{\small\color{refgray}数据：第二轮人员优化涉及约4000名员工\par}
{\small\color{refgray}边际变化：市场此前关注销量下滑，GS强调911组合定价和成本优化足以驱动利润改善，第二轮人员优化规模超预期。\par}
\subsubsection*{关键数据}
\begin{itemize}
  \item 耐克Lifestyle占品牌收入约60\%，FY27仍中个位数下滑
  \item 东京百货珠宝销售2Q26同比增20\%以上
  \item LVMH时装皮具2Q26有机销售增速预期+3\%，高于市场共识+1.7\%
  \item B站搜索广告收入同比+100\%，播放页广告+80\%
  \item 宝马中国区销量同比-30\%，欧洲+8\%，美国+10\%
  \item 保时捷2Q26利润率8.5\%，市场共识6.6\%
  \item 古茗2026年咖啡GMV约80亿元，为瑞幸低双位数百分比
  \item 茶百道6月同店销售下滑低个位数
\end{itemize}
\begin{figure}[htbp]
\centering
\includegraphics[width=0.88\linewidth]{figures/fig_016.jpg}
\caption{Exhibit 1 - EXHIBIT 3: In FY27, we expect Performance categories to continue gaining share with the sales mix Nike Brand Revenue Split, 2027E}
\end{figure}
\begin{figure}[htbp]
\centering
\includegraphics[width=0.88\linewidth]{figures/fig_024.jpg}
\caption{EXHIBIT 1 - EXHIBIT 2: Our weighted-average model predicts a +12.1\% OSG for Richemont Group in 1Q27, c. +110bps above company-collected consensus, supported by jewellery momentum in Japan}
\end{figure}
\begin{figure}[htbp]
\centering
\includegraphics[width=0.88\linewidth]{figures/fig_042.jpg}
\caption{Exhibit 1 - Exhibit 1: Bili maintained solid user engagement growth into early 2026 Exhibit 2: Ads growth is accelerating, consistent with timespent trend}
\end{figure}
\begin{figure}[htbp]
\centering
\includegraphics[width=0.88\linewidth]{figures/fig_049.jpg}
\caption{Exhibit 3 - Exhibit 3: BMW's YTD price performance is primarily explained by multiple contraction and, following the June 19 profit warning, EPS downgrades. BMW 26YTD price performance decomposition Exhibit 4: BMW is currently trading close}
\end{figure}
{\small\color{refgray}本节综合 6 篇研究；涉及 Bernstein、GS。}
\newpage
\subsection{区域市场与主题：代币化资产落地、中国准财政工具与太空竞赛新格局}
\textbf{本期区域市场与主题聚焦三大结构性变化：代币化真实世界资产（RWA）从实验进入部署期，Robinhood Chain数据验证用户对代币化资产流动性的需求；中国通过新型政策性资金工具（NPBFI）绕过地方债务约束，定向支持高技术转型；太空领域中国长征十号回收成功加速追赶，但SpaceX在复用深度和规模化上仍保持绝对领先。}
\subsubsection*{主流共识}
\begin{itemize}
  \item 代币化RWA市场正从发行实验转向效用落地，Robinhood Chain等实际部署验证了用户对代币化资产流动性的需求。
  \item 中国NPBFI工具是当前双重约束（支持高技术转型与控制地方债务）下的理性准财政替代方案，资金集中于12个经济大省的战略领域。
  \item 中国已成为SpaceX在航天发射领域最主要的长期竞争者，长征十号回收成功比预期提前约半年。
\end{itemize}
\subsubsection*{分歧与条件差异}
\begin{itemize}
  \item 代币化RWA的驱动力：Bernstein强调Robinhood Chain的DEX交易量和公司代币持有者增长是信号，而市场部分观点仍视其为Meme币驱动的短期现象。
  \item 中国NPBFI对GDP的拉动效果：GS基准情景估算拉动0.5个百分点，但效果集中在2026年底至2027年初，市场对时点和幅度存在分歧。
  \item 太空竞争格局：Bernstein认为中国在加速追赶但差距仍大，MS则通过SpaceX入列Space 60指数强调其已从初创公司变为运营基准，两者对追赶速度的判断有差异。
\end{itemize}
\subsubsection*{机构观点}
\paragraph{Bernstein} 代币化RWA市场市值已突破510亿美元，年内增长50\%，资产持有者突破100万，增长75\%。Robinhood Chain上线不到15天，DEX交易量达31亿美元，跻身全网前五；链上已有6.5万用户持有1300万美元公司代币和3亿美元稳定币。公司代币覆盖90余只标的，支持24x7交易并可作为借贷抵押品。报告认为，代币化正从发行实验转向效用落地，Robinhood Chain证明用户愿意为代币化资产的流动性买单。
{\small\color{refgray}数据：代币化RWA市值510亿美元，年内增长50\%\par}
{\small\color{refgray}数据：资产持有者突破100万，年内增长75\%\par}
{\small\color{refgray}数据：Robinhood Chain DEX 7天交易量31亿美元，全网前五\par}
{\small\color{refgray}数据：6.5万用户持有1300万美元公司代币\par}
{\small\color{refgray}边际变化：此前市场关注代币化能做什么，现在Robinhood Chain的实际部署开始回答代币化资产怎么用、谁在用、用在哪。\par}
\paragraph{GS} 中国2026年NPBFI规模扩至8000亿元，通过三大政策银行以股权形式投向国家发改委筛选的战略项目，填补15\%-25\%的资本金缺口。资金成本低（PSL利率1.75\%，国开债收益率1.5\%-1.6\%），80\%流向12个经济大省（贡献69\%GDP和80\%研发支出），集中于AI基础设施、半导体供应链、绿色转型及六张网项目。GS估算基准情景下拉动实际GDP约0.5个百分点，效果集中在2026年底至2027年初。
{\small\color{refgray}数据：NPBFI规模8000亿元\par}
{\small\color{refgray}数据：PSL利率1.75\%，国开债收益率1.5\%-1.6\%\par}
{\small\color{refgray}数据：80\%资金流向12个经济大省\par}
{\small\color{refgray}数据：12省贡献69\%GDP和80\%研发支出\par}
{\small\color{refgray}边际变化：NPBFI从2022年疫情期间的临时工具（7400亿元）制度化，2026年规模扩至8000亿元，成为中央统筹、不增加地方隐性债务的准财政替代方案。\par}
\paragraph{Bernstein} 中国长征十号B火箭一级助推器回收成功，比预期提前约半年，但仅具备一级复用能力，每次发射仍需新造二级火箭。SpaceX猎鹰九号已实现近十年复用，部分助推器复用超20次，2025年完成165次发射。星舰全复用设计可将年发射量提升一个数量级，Bernstein预测2031年SpaceX年发射量约3500次，长征十号届时可能难以突破100次。中国已向国际电联申报超20万颗低轨卫星星座计划，但更多是锁频段策略。
{\small\color{refgray}数据：长征十号B一级回收成功，比预期提前约半年\par}
{\small\color{refgray}数据：SpaceX猎鹰九号部分助推器复用超20次\par}
{\small\color{refgray}数据：2025年猎鹰九号完成165次发射\par}
{\small\color{refgray}数据：星舰全复用，长征十号仅一级复用\par}
{\small\color{refgray}边际变化：中国航天回收能力从零到一，但全复用才是成本指数级下降的关键，SpaceX在复用深度和规模化上仍绝对领先。\par}
\paragraph{MS} MS将SpaceX纳入Space 60指数，描述为'X of 1'企业，横跨发射、卫星通信和AI基础设施。截至2026年3月，SpaceX完成约650次轨道发射，99\%成功率，Starlink超10,000颗在轨卫星（占全球可机动卫星75\%），服务约1200万宽带用户。星舰有效载荷是猎鹰9号四倍以上，预计2030年发射成本降至约500美元/公斤，2040年降至150美元/公斤以下。报告预测2026年营收450亿美元，2040年达3.3万亿美元，但对应每年约3000亿美元资本支出存在不确定性。
{\small\color{refgray}数据：SpaceX完成约650次轨道发射，99\%成功率\par}
{\small\color{refgray}数据：Starlink超10,000颗在轨卫星，占全球可机动卫星75\%\par}
{\small\color{refgray}数据：宽带用户约1200万\par}
{\small\color{refgray}数据：星舰有效载荷是猎鹰9号四倍以上\par}
{\small\color{refgray}边际变化：太空产业链从少数玩家主导的发射竞赛转向有分层、有现金流、有退出路径的成熟行业结构，SpaceX入列指数标志其成为太空基础设施的运营基准。\par}
\paragraph{NOM} 市场对沪电股份的担忧（英伟达Kyber机架延迟、客户分散供应链）可能过度。真正驱动价值的是2026年下半年Rubin GPU量产、ASIC客户进入及光模块PCB放量。HDI产能仍是战略资产，技术壁垒和材料设备短缺限制竞争者进入。预计HDI业务2026-2028财年收入年复合增长率82\%，到2028财年贡献总收入63\%。ASIC客户从2028年起可能采用HDI，进一步拉长需求周期。
{\small\color{refgray}数据：HDI业务2026-2028财年收入年复合增长率82\%\par}
{\small\color{refgray}数据：2028财年HDI贡献总收入63\%\par}
{\small\color{refgray}数据：Rubin GPU PCB出货2026年中开始\par}
{\small\color{refgray}数据：ASIC客户2028年起可能采用HDI\par}
{\small\color{refgray}边际变化：市场聚焦单一产品（Kyber）延迟，但Rubin GPU的HDI价值含量更高，价值升级路径未中断，只是节奏切换。\par}
\subsubsection*{关键数据}
\begin{itemize}
  \item 代币化RWA市值510亿美元，年内增长50\%
  \item Robinhood Chain DEX 7天交易量31亿美元，全网前五
  \item 中国NPBFI规模8000亿元，PSL利率1.75\%
  \item SpaceX完成约650次轨道发射，99\%成功率
  \item Starlink超10,000颗在轨卫星，占全球可机动卫星75\%
  \item 长征十号B一级回收成功，比预期提前约半年
  \item HDI业务2026-2028财年收入年复合增长率82\%
\end{itemize}
\begin{figure}[htbp]
\centering
\includegraphics[width=0.88\linewidth]{figures/fig_006.jpg}
\caption{EXHIBIT 1 - EXHIBIT 2: RWA holders on Robinhood Chain have surpassed 67K users Robinhood Chain - RWA Holders (in '000s) EXHIBIT 3: DEX volumes have crossed \textbackslash{}\$3Bn in last 7 days}
\end{figure}
\begin{figure}[htbp]
\centering
\includegraphics[width=0.88\linewidth]{figures/fig_029.jpg}
\caption{Exhibit 2 - Exhibit 2: How the NPBFI mechanism works \#\# Why have Chinese policymakers increasingly favored this tool in recent years? In our view, China's fiscal system now faces a dual mandate: supporting the economy's transition toward hig}
\end{figure}
\begin{figure}[htbp]
\centering
\includegraphics[width=0.88\linewidth]{figures/fig_039.jpg}
\caption{Exhibit 2 - EXHIBIT 1: Long March 10B cable system catch EXHIBIT 2: SpaceX chopstick catch \#\# COMPETITIVE PICTURE}
\end{figure}
\begin{figure}[htbp]
\centering
\includegraphics[width=0.88\linewidth]{figures/fig_092.jpg}
\caption{Exhibit 1 - Exhibit 1: The MS Space 60 \#\# For further reading: • Space: The Space 60: Picks \& Shovels for the Final Frontier (12 Apr 2026) \textbackslash{}- SpaceX: AI's Final Frontier; Initiate at Overweight, PT \textbackslash{}\$300 (7 Jul 2026)}
\end{figure}
{\small\color{refgray}本节综合 5 篇研究；涉及 Bernstein、GS、NOM、MS。}
\newpage
\section{辅助信号｜战略咨询}
本板块整合波士顿咨询两份报告，聚焦两个主题：一是泰国AI基建的窗口期与治理挑战，二是国防飞机维护系统的效率瓶颈。两份报告均指向资源投入与产出之间的转化效率问题，强调流程、治理和时机对战略结果的决定性影响。
\subsection{泰国AI基建：窗口期与治理挑战}
\textbf{泰国成为东南亚AI枢纽的窗口期正在关闭，关键在于治理效率而非资源禀赋。全球数据中心供需收敛前，泰国需在3年内建立跨部门协调机制、前置基础设施投资和明确的数据治理规则，否则将被马来西亚等先行者锁定供给。}
\begin{itemize}
  \item 全球数据中心短缺带来的运营商议价优势不会持续太久，供需收敛正在加速，窗口期约3年。
  \item 泰国拥有电价、土地、地缘中立等结构性优势，2025年研究促进委员会收到230亿美元数据中心研究申请，2026年字节跳动承诺投入250亿美元。
  \item 马来西亚案例显示，公共基础设施先行（TNB投入188亿林吉特扩建输电网，输电可靠性提升65\%）可将电网接入时间从3-4年压缩至12个月，泰国尚未建立类似机制。
  \item 泰国目前23个部委参与数据中心审批，缺乏跨部门协调机制，治理效率是最大瓶颈。
  \item 数据治理确定性是长期基础设施承诺的前提，需明确数据驻留规则、充分性认定和跨境传输机制。
\end{itemize}
\begin{figure}[htbp]
\centering
\includegraphics[width=0.88\linewidth]{figures/fig_136.jpg}
\caption{Exhibit 1 - Delivering on that ambition requires the right infrastructure. For AI, that infrastructure is the data center. (Exhibit 1.) AI-era data centers are not simply larger versions of what came before. Higher power densities, purpose-built cooling systems, and netwo}
\end{figure}
\begin{figure}[htbp]
\centering
\includegraphics[width=0.88\linewidth]{figures/fig_139.jpg}
\caption{EXHIBIT 6 - Global foundation model training is the hardest layer to compete for. These workloads are largely location-agnostic and will follow wherever infrastructure quality and market conditions are most favorable. \#\# EXHIBIT 6 Thailand needs to build an additional 1.6}
\end{figure}
\subsection{国防飞机维护：资源投入与出勤率的背离}
\textbf{美国及盟国军机任务可用率持续下降，问题不在资源总量，而在将资源转化为出勤结果的维护系统效率。通过改进流程、优先级排序和治理机制，可在6-12个月内将可用率提升30\%-50\%，无需额外投入。}
\begin{itemize}
  \item 美国空军2024年全机队任务可用率降至67\%，为至少20年来最低；英国F-35机队2024年仅三分之一完全可用。
  \item 美国作战与维护支出经通胀调整后约为2000年的两倍，但空军和海军航空兵可用率同期下降12\%-19\%。
  \item 法国战斗机与战术运输机可用率较2014年分别下降26.5和15个百分点；德国战斗机平均可用率仅64\%，远低于80\%基准。
  \item 制约因素包括：维修时间差异大（同型号飞机不同基地检查时间差2-3倍）、缺乏全局优先级排序、短期飞行计划压倒长期战备。
  \item 通过直接改进维护流程、需求纪律和治理机制，可在6-12个月内将任务可用率提升30\%-50\%，且无需额外投入。
\end{itemize}
\begin{figure}[htbp]
\centering
\includegraphics[width=0.88\linewidth]{figures/fig_136.jpg}
\caption{Exhibit 1 - Delivering on that ambition requires the right infrastructure. For AI, that infrastructure is the data center. (Exhibit 1.) AI-era data centers are not simply larger versions of what came before. Higher power densities, purpose-built cooling systems, and netwo}
\end{figure}
\newpage
\section{辅助信号｜智库/国际机构}
本板块聚焦国际机构最新研究，涵盖数字支付基础设施、劳动力市场结构性挑战及财政可持续性三大主题。亚开行报告揭示尼泊尔支付系统从互操作走向统一标准的质变，以及中国ICT监管改革的贸易乘数效应；BIS与IMF报告分别指出零售支付竞争格局的结构性壁垒和BigTech金融服务的监管盲区；OECD与IMF则从劳动力市场政策评估、竞业限制、AI影响及欧洲财政紧缩等角度，提供了多维度政策启示。
\subsection{数字支付与金融科技：基础设施、竞争与监管}
\textbf{数字支付正经历从技术采纳到基础设施统一的关键转折，但后端清算网络的高集中度与BigTech的跨业务风险，要求监管框架从单一机构监管转向系统性、互操作性导向的治理。}
\begin{itemize}
  \item 亚开行报告指出，尼泊尔央行2025年发布的国家支付交换机蓝图，标志着从分散系统向共享数字公共基础设施的转变，核心是强制实施开放API规则手册和NEPALPAY二维码标准，以解决多系统碎片化问题，并为跨境支付铺路。
  \item BIS公告显示，全球两大卡组织Visa和Mastercard在主要地区的合计市场份额稳定在约95\%，利润率持续高位，说明后端清算网络的高固定成本和强网络效应构成天然垄断壁垒，前端应用层创新未改变竞争结构。
  \item IMF技术说明强调，BigTech在新兴市场通过支付和信贷业务快速规模化，如肯尼亚M-Pesa控制近99\%移动支付份额，且支付数据被用于信贷决策，形成业务互联风险；现有监管框架无法有效覆盖其跨业务、跨市场的集团化运作。
  \item 亚开行关于中国数字服务贸易的论文量化显示，中国ICT行业单边监管改革（如提高监管问责、建立竞争框架）的贸易拉动效应远超RCEP协定实施，改革重点在于跨境数据流动和数据本地化要求。
  \item BIS报告区分了基于卡的网络和基于账户对账户的快速支付系统（FPS），指出FPS在新兴市场改变游戏规则，但前端竞争若不伴随互操作性，可能导致碎片化，损害规模效率。
\end{itemize}
\begin{figure}[htbp]
\centering
\includegraphics[width=0.88\linewidth]{figures/fig_097.jpg}
\caption{Figure 1 - Figure 1: High-Level Architecture of the Retail Payment Switch ACH = automated clearing house, API = application programming interface, BFIs = banks and financial institutions, Co. = company, GL = general ledger, Govt = governme}
\end{figure}
\begin{figure}[htbp]
\centering
\includegraphics[width=0.88\linewidth]{figures/fig_105.jpg}
\caption{Figure 1 - \textbackslash{}- Aggregator SuperApps provide access to various third-party providers on their service, bringing together various products and services into one easy-to-access application. Examples include Gojek and Grab. It is important to note that the distinction between}
\end{figure}
\subsection{劳动力市场结构性挑战：AI冲击、竞业限制与非正规就业}
\textbf{劳动力市场面临AI替代、竞业限制泛化及非正规就业停滞等多重结构性压力，政策需从单一就业率目标转向兼顾社会效果与合规成本降低。}
\begin{itemize}
  \item IMF对以色列的AI影响评估显示，约三分之二工人面临高AI暴露，其中三分之一处于“高暴露、低互补”区间（最易被替代），集中在ICT、金融和保险行业；公共部门因教育、医疗等职业AI互补性强，替代风险更低。
  \item OECD加拿大竞业限制调查发现，29\%-39\%私营部门工人受约束，且低薪、固定期限合同工同样频繁出现；超过一半受约束工人同时被施加三项以上限制条款，许多条款在法律上脆弱（无明确期限、地域或补偿），但工人因声誉担忧和合同误判仍选择遵守。
  \item OECD哥伦比亚劳动市场评估指出，2025年综合劳动改革虽提升合规标准，但非正规就业率自2024年中停止下降，仍达57\%（农村84\%），因改革增加带薪休假、加班补偿等条款推高正式用工成本，企业有动力规避合规。
  \item OECD积极劳动力市场政策（ALMP）报告显示，37个国家已实施200多项嵌入社会目标（如增强求职者主体感、心理健康）的政策，但仅约十分之一国家系统评估社会效果，评估体系滞后于政策设计。
  \item IMF以色列报告揭示一个反直觉发现：高学历和年轻群体反而更脆弱，因其集中在ICT等AI高暴露行业，而低学历群体多从事体力劳动，AI替代风险相对较低。
\end{itemize}
\begin{figure}[htbp]
\centering
\includegraphics[width=0.88\linewidth]{figures/fig_116.jpg}
\caption{Figure 1 - \#\# B. Preparedness and Challenges 4. Israel is relatively well prepared for AI adoption. The IMF's AI Preparedness Index ranks Israel 18th globally, reflecting strong innovation capacity and a supportive environment, while the Stanford AI Index places Israel a}
\end{figure}
\begin{figure}[htbp]
\centering
\includegraphics[width=0.88\linewidth]{figures/fig_126.jpg}
\caption{Figure 1 - \#\# Figure 1. Non-compete and related clauses are relatively common in Canada Share of workers bound by each employment contract clause A comparison of workers with and without non-compete clauses shows that these agreements are present across all types of empl}
\end{figure}
\subsection{财政可持续性与外部冲击：欧洲紧缩与埃塞俄比亚改革}
\textbf{欧洲面临老龄化、能源转型和国防三重支出压力，需结构性改革与财政整顿并行；埃塞俄比亚在外部战争冲击下坚持汇率市场化改革，IMF项目框架提供政策取舍的真实样本。}
\begin{itemize}
  \item IMF欧洲财政报告量化显示，到2040年，老龄化、能源转型和国防三项刚性支出将推动欧洲平均政府支出增加约5个百分点GDP；若不行动，平均债务率将在15年内翻倍至130\%GDP。
  \item IMF模拟表明，中等力度改革组合（产品市场、劳动力市场、养老金调整等）可填补约三分之一财政缺口，但改革效果滞后于支出压力释放，需中期财政整顿作为第二支柱。
  \item IMF埃塞俄比亚第五次审议显示，中东战争通过三条渠道传导：燃料进口成本飙升（柴油现货价一度达270美元/桶）、化肥采购不确定性、黄金出口市场中断（依赖包机运往迪拜精炼）。
  \item 埃塞俄比亚在外部冲击下仍坚持汇率市场化改革，所有量化绩效指标均达标，2024/25财年实际GDP增长9.2\%，通胀从26.6\%降至11.7\%，外汇储备改善，显示改革韧性。
  \item IMF报告强调，欧洲养老金改革和促进增长的国内改革在缓解财政压力方面效果最突出，但时间错配意味着改革需与其他工具（如重新定义政府边界）配合，才能避免债务失控。
\end{itemize}
\begin{figure}[htbp]
\centering
\includegraphics[width=0.88\linewidth]{figures/fig_106.jpg}
\caption{Figure 1 - Many European countries find themselves in a difficult budgetary position. They are facing new shocks, even as the legacy of past shocks still looms large, with public debt having increased sharply during the COVID-19 pandemic and remaining high. As financial }
\end{figure}
\begin{figure}[htbp]
\centering
\includegraphics[width=0.88\linewidth]{figures/fig_111.jpg}
\caption{Figure 1 - restrictions. The restrictions currently in place for BOP reasons are the imposition of hard ceilings on access to and use of FX for travel purposes, and the imposition of hard ceilings on access to and use of FX for the purposes of cross-border payment of mod}
\end{figure}
\section{结语}
后续需验证的关键节点包括：美国电力需求增速的实际落地情况（数据中心建设进度与电气化渗透率）、中国互联网板块盈利拐点是否如期在2Q-3Q26出现、代币化RWA的流动性需求能否持续转化为实际交易量、以及中国NPBFI对GDP的拉动效果是否在预期时点兑现。
\newpage
\section{覆盖概览}
投行/券商 30 篇；战略咨询 2 篇；智库/国际机构 10 篇。\par
投行覆盖：GS 15篇 / Bernstein 7篇 / JPM 3篇 / MS 3篇 / NOM 2篇\par
\section{Disclaimer}
{\scriptsize\color{refgray}
This community is a paid learning and information-sharing community created for general educational purposes only. The membership fee is charged solely for access to community discussions, educational content, organization, and operational costs. It is not an advisory fee, management fee, performance fee, brokerage fee, or compensation for personalized financial, investment, legal, tax, or professional advice.\par
I participate and share information solely as an individual and for educational discussion among community members. I am not acting as a financial adviser, investment adviser, broker, dealer, portfolio manager, legal adviser, tax adviser, accountant, fiduciary, or any other licensed professional adviser. Nothing shared in this community constitutes, and should not be interpreted as, financial advice, investment advice, legal advice, tax advice, a recommendation, solicitation, offer, endorsement, instruction, or guarantee to buy, sell, hold, short, trade, or invest in any security, cryptocurrency, fund, derivative, strategy, or other financial product.\par
All content shared in this community, including messages, discussions, links, files, charts, examples, market commentary, watchlists, AI-generated or AI-polished summaries, and third-party materials, is general, impersonal, and educational in nature. It does not take into account any individual's financial situation, investment objectives, risk tolerance, investment horizon, liquidity needs, tax status, legal restrictions, or suitability requirements. No content should be relied upon as suitable for any specific person, account, portfolio, or financial decision.\par
Information shared in this community may be incomplete, inaccurate, outdated, speculative, AI-generated, AI-edited, or based on third-party internet sources that have not been independently verified. I make no representation or warranty as to the accuracy, completeness, timeliness, reliability, or suitability of any information shared.\par
Financial markets involve substantial risk, including the possible loss of principal. Past performance, historical data, hypothetical examples, simulated results, backtests, personal opinions, or educational examples do not guarantee future results. Each member is solely responsible for conducting their own research, exercising independent judgment, and making their own decisions. Before making any financial, investment, legal, or tax decision, members should consult qualified and licensed professionals.\par
Participation in this community, payment of any membership fee, or communication with me or other members does not create any adviser-client, fiduciary, professional, contractual, agency, partnership, employment, or other advisory relationship. By joining or participating in this community, you acknowledge and agree that you are solely responsible for your own decisions, actions, transactions, profits, losses, risks, and consequences, and that I shall not be liable for any direct, indirect, incidental, consequential, financial, legal, tax, or other loss or damage arising from or related to any content, discussion, information, or material shared in this community.\par
}
\newpage
\begin{center}
{\Large 更多详情报告kcdesk.com}\par\vspace{0.8cm}
\includegraphics[width=0.58\linewidth]{\detokenize{../../prompts/zsxq_img.jpg}}
\end{center}
\end{document}
